\begin{center}
\begin{tikzpicture}[x=.88cm,y=.88cm,font=\small]
  \fill[paper] (0,0) rectangle (12.3,6.65);

  \draw[source!70,line width=.8pt] (.65,2.1) .. controls (3.3,1.7) and
    (6.1,2.5) .. (11.7,1.78);
  \node[text=source,below,font=\scriptsize] at (7.8,1.93) {长江方向（示意）};

  \fill[dynasty!23] (.9,2.75) -- (2.18,2.08) -- (3.48,2.52) --
    (3.55,3.78) -- (2.36,4.36) -- (1.08,3.88) -- cycle;
  \node[align=center] at (2.22,3.2) {关中\\长安};

  \fill[timeline!26] (3.95,4.1) ellipse (1.18 and .82);
  \node[align=center] at (3.95,4.1) {平阳\\汉（赵）中心};

  \fill[dynasty!39] (5.78,3.06) ellipse (1.22 and .86);
  \node[align=center,text=white] at (5.78,3.06) {洛阳\\西晋京师};

  \fill[source!24] (8.28,4.55) ellipse (1.24 and .8);
  \node[align=center] at (8.28,4.55) {并、冀一带\\起兵与军镇};

  \fill[timeline!22] (10.42,1.05) ellipse (1.34 and .7);
  \node[align=center] at (10.42,1.05) {建业（建康）\\江东军政重组};

  \fill[source!13] (7.32,1.06) ellipse (1.1 and .57);
  \node[align=center,font=\scriptsize] at (7.32,1.06) {淮南—江北\\渡江通道};

  \draw[->,very thick,color=dynasty]
    (5.43,3.52) .. controls (4.92,3.94) and (4.66,4.08) .. (4.52,4.1);
  \node[text=dynasty,above,font=\scriptsize,align=center] at (4.94,4.72)
    {311年洛阳陷落后\\怀帝被押往平阳};

  \draw[->,very thick,color=dynasty]
    (3.34,2.72) .. controls (3.72,3.09) and (4.13,3.37) .. (4.58,3.64);
  \node[text=dynasty,below,font=\scriptsize,align=center] at (3.77,2.43)
    {316年长安失守\\北方朝廷终结};

  \draw[->,thick,dashed,color=timeline]
    (6.75,2.64) .. controls (7.72,2.04) and (8.78,1.41) .. (9.28,1.18);
  \node[text=timeline,above,font=\scriptsize,align=center] at (8.37,2.18)
    {307年司马睿出镇建业；\\南下与流民渡江并非单一路线};

  \draw[->,thick,dashed,color=source]
    (8.12,3.8) .. controls (7.5,3.38) and (6.99,3.15) .. (6.92,3.12);
  \node[text=source,right,font=\scriptsize] at (7.12,3.72) {战乱压力};

  \node[draw=source!75,fill=paper,rounded corners=2pt,align=center,
    inner sep=3pt,font=\scriptsize] at (8.85,5.72)
    {永嘉之乱不是一次单向撤离：\\京师崩解、北方政权兴起与江东重组同时展开};
\end{tikzpicture}

\smallskip
\small\textit{图\quad 西晋崩解与永嘉前后的几条政治、人员移动方向，约307--318年：据《晋书·惠帝纪》《怀帝纪》《元帝纪》及《刘元海载记》，并参照《资治通鉴》卷八十五至八十九，概括洛阳、平阳、长安与建业之间的政局关联。节点方位约略，色块只标示政治和军事重心，绝非疆界；箭线只提示史料可见的方向与关联，不复原个人南渡、俘押或军队行动的实际路线、规模和时序。}
\end{center}
